\chapter{Web Server Interface and Cloud Integration}

This chapter refines the web server and cloud integration so the report is accurate and presentation-ready. The ESP32 runs AP and STA together: the AP keeps local access alive for control/OTA, while STA reconnects in the background to restore cloud links and telemetry.

\section{Web Server Design}

Static UI assets live in LittleFS and are served by \texttt{ESPAsyncWebServer} so request handling stays asynchronous and non-blocking. Figure~\ref{fig:backend_block} summarizes how the networking stack, web layer, message queues, and FreeRTOS tasks interact from ingress to actuation/cloud forwarding.

\begin{figure}[H]
    \centering
    % TODO: Thay bằng sơ đồ khối backend thực tế
    \fbox{\parbox{0.9\textwidth}{\centering
        \vspace{2.0cm}
        \textit{[Backend Block Diagram]}\\
        AP/STA WiFi $\rightarrow$ ESPAsyncWebServer + WebSocket $\rightarrow$ Message Queue $\rightarrow$ FreeRTOS Tasks (Sensor, Control, Cloud) $\rightarrow$ Peripherals/Cloud\\
        \vspace{2.0cm}
    }}
    \caption{Sơ đồ khối backend: từ tầng mạng đến queue và các task điều khiển/đồng bộ cloud.}
    \label{fig:backend_block}
\end{figure}

\subsection{Access Point Configuration}

The ESP32 always boots in AP+STA mode so it never becomes unreachable. The AP is fixed at SSID \texttt{YoloUNO\_AP}, static IP \texttt{192.168.4.1}, WPA2-PSK. In \texttt{task\_wifi.cpp}, \texttt{ensureAP()} keeps that AP alive while \texttt{startSTA()} tries the stored home WiFi; if the STA link drops, \texttt{Wifi\_reconnect()} retries the uplink but still preserves the AP. Practically, users can always reach \texttt{192.168.4.1} for local control/OTA even when home WiFi changes; once STA succeeds, the device also announces its LAN IP and OTA URL.

\subsection{Request Handling and Task Flow}

All traffic enters the asynchronous web layer in \texttt{task\_webserver.cpp}. Static files, REST endpoints, and WebSocket events share a consistent path: parse/validate input, execute the requested control or info action, then broadcast the refreshed state to connected clients. Key endpoints include health check (\texttt{/health}), device info (\texttt{/info}), CSV lifecycle (\texttt{/download}, \texttt{/csv-info}, \texttt{/clear}), and OTA hooks (ElegantOTA). WebSocket messages are handled in \texttt{onEvent()} and forwarded to \texttt{handleWebSocketMessage()} for device control.

\textbf{Backend execution.} \texttt{connnectWSV()} (in \texttt{task\_webserver.cpp}) mounts LittleFS assets, registers REST/WebSocket handlers, and enables ElegantOTA before \texttt{server.begin()}. \texttt{onEvent()} parses WebSocket frames and hands them to \texttt{handleWebSocketMessage()} for LED/pump control and state fan-out. \texttt{task\_wifi.cpp} enforces AP+STA so the AP at \texttt{192.168.4.1} stays up during STA reconnects. In the cloud path, \texttt{CORE\_IOT\_task} keeps the MQTT session alive, publishes telemetry every 10s, and applies RPC commands (fan/LED/pump) while mirroring states back to keep the dashboard switches in sync. Health, info, CSV download/clear/info, and OTA progress are exposed via lightweight endpoints and WebSocket notifications.
\textbf{RPC and reconnect.} RPC callbacks (\texttt{setFanStatus}, \texttt{setLedEnabled}, \texttt{setWaterPumpStatus}) are executed inside \texttt{CORE\_IOT\_task}; each command drives the actuator, updates shared sensor state, and the paired getter RPCs return the latest value so dashboard toggles stay aligned. The MQTT loop reconnects whenever the link drops, resubscribes RPC/shared attributes, and resumes publishing without losing state. On the WiFi side, \texttt{Wifi\_reconnect()} rate-limits STA retries every 30s but never drops the AP, ensuring local UI/WebSocket/OTA remain reachable during WAN outages.

\begin{figure}[H]
    \centering
    % TODO: Replace with actual backend flowchart
    \fbox{\parbox{0.9\textwidth}{\centering
        \vspace{1.8cm}
        \textit{[Backend Flowchart]}\\
        Client HTTP/WebSocket $\rightarrow$ Async Handler $\rightarrow$ Validate/Dispatch $\rightarrow$ Control or Info Action $\rightarrow$ State Broadcast\\
        \vspace{1.8cm}
    }}
    \caption{Request/response flow across HTTP/WebSocket handlers, control actions, and outbound broadcasts.}
    \label{fig:backend_flow}
\end{figure}

\subsection{User Interface Implementation}

UI assets (HTML5/CSS3/JavaScript) are decoupled from firmware—update LittleFS only to change the look. WebSocket pushes keep WiFi/MQTT badges live without reloading the page.

\begin{enumerate}
    \item \textbf{Device Control Panel:} LED (GPIO 48) and Pump/Fan (GPIO 2) toggles.
    \item \textbf{Real-time Status:} Temperature, humidity, light level, WiFi/cloud indicators.
    \item \textbf{Cloud Sync State:} MQTT connection badge and last publish timestamp.
\end{enumerate}

\begin{figure}[H]
    \centering
    % TODO: Replace placeholder image paths with real screenshots
    \begin{minipage}{0.45\textwidth}
        \centering
        \includegraphics[width=0.9\textwidth]{assets/ui_home_placeholder.png}
        \caption*{Home — overview + quick toggles}
    \end{minipage}
    \hfill
    \begin{minipage}{0.45\textwidth}
        \centering
        \includegraphics[width=0.9\textwidth]{assets/ui_devices_placeholder.png}
        \caption*{Devices — LED/Pump controls + live metrics}
    \end{minipage}

    \vspace{0.5cm}

    \begin{minipage}{0.45\textwidth}
        \centering
        \includegraphics[width=0.9\textwidth]{assets/ui_info_placeholder.png}
        \caption*{Info — device ID, firmware, IP, uptime}
    \end{minipage}
    \hfill
    \begin{minipage}{0.45\textwidth}
        \centering
        \includegraphics[width=0.9\textwidth]{assets/ui_settings_placeholder.png}
        \caption*{Settings — WiFi credentials, MQTT token, publish interval}
    \end{minipage}
    \caption{Placeholders for UI screens: Home, Devices, Info, Settings.}
    \label{fig:web_ui_screens}
\end{figure}

\subsection{Web Communication Protocols}

\subsubsection{HTTP vs. WebSocket}
Traditional web servers use the HTTP protocol, which follows a request-response model. The client requests data, the server responds, and the connection is closed. This is inefficient for real-time control because the client must constantly poll the server to get updates.

For our control interface, we chose the \textbf{WebSocket} protocol. WebSocket provides a full-duplex communication channel over a single TCP connection. Once established, the server can push data to the client immediately without waiting for a request. This significantly reduces latency and network overhead, making it ideal for controlling actuators (like LEDs) and displaying live sensor data.

\subsection{WebSocket Communication}

To achieve real-time performance without page reloads, we utilized the WebSocket protocol. Unlike traditional HTTP polling, WebSocket establishes a persistent, full-duplex connection between the browser and the ESP32.

\begin{lstlisting}[language=C++, caption={WebSocket event handler for real-time control}]
void onEvent(AsyncWebSocket *server, AsyncWebSocketClient *client, 
             AwsEventType type, void *arg, uint8_t *data, size_t len) {
    if (type == WS_EVT_DATA) {
        AwsFrameInfo *info = (AwsFrameInfo *)arg;
        if (info->opcode == WS_TEXT) {
            String message = String((char *)data).substring(0, len);
            
            // Parse JSON command
            // Example: {"action":"toggleLed", "value":true}
            handleWebSocketMessage(message);
        }
    }
}
\end{lstlisting}

When a user clicks a button on the web interface, a JSON message is sent over the WebSocket. The ESP32 parses this message and triggers the corresponding FreeRTOS task (e.g., giving a semaphore to the LED task). Conversely, when sensor values change, the ESP32 broadcasts the new data to all connected clients immediately.

\begin{figure}[H]
    \centering
    % TODO: Replace with actual diagram of WebSocket flow
    \fbox{\parbox{0.8\textwidth}{\centering
        \vspace{1.5cm}
        \textit{[Insert diagram of WebSocket Communication Flow here]}\\
        \textit{Client -> JSON Command -> ESP32 -> Task Notification}\\
        \vspace{1.5cm}
    }}
    \caption{WebSocket communication flow for real-time device control.}
    \label{fig:websocket_flow}
\end{figure}

\section{Cloud Integration and Data Publishing}

While the local web server provides immediate control, long-term data storage and remote access are handled by the CoreIOT cloud platform.

\subsection{Station Mode Connectivity}

The device connects to the local WiFi network in Station (STA) mode. We implemented a robust reconnection state machine that:
\begin{enumerate}
    \item Attempts to connect to the configured WiFi SSID.
    \item If connection fails, it waits for a backoff period before retrying.
    \item Periodically checks connection health and triggers reconnection if the link is dropped.
    \item Synchronizes system time via NTP (Network Time Protocol) upon successful connection.
\end{enumerate}

Crucially, the reconnection logic runs in a low-priority background task to avoid blocking critical sensor processing.
\subsection{Firmware Updates (OTA)}
The system supports Over-The-Air (OTA) firmware updates, allowing for convenient maintenance without physical access to the device.
\textbf{Initial Setup:} The first firmware upload must be performed via the USB-C interface to install the bootloader and the initial application containing the OTA logic.
\textbf{Subsequent Updates:} Once the initial firmware is running, subsequent updates can be performed wirelessly. During this process, the device can remain powered by the external 12V power supply, eliminating the need to reconnect the USB-C cable. This feature is particularly valuable for deployed devices where the USB port may be inaccessible.
\subsection{Data Publishing to CoreIOT}

We utilize the MQTT protocol to publish sensor telemetry to the CoreIOT server (\texttt{app.coreiot.io}). The publishing logic is encapsulated in the \texttt{CORE\_IOT\_task}, which aggregates data from the shared sensor queue and transmits it securely.

\textbf{Authentication Configuration:}
The connection is authenticated using a unique Device Token. This token acts as the MQTT username, ensuring that data is routed to the correct device dashboard in the cloud.

\begin{lstlisting}[language=C++, caption={CoreIOT authentication and data publishing}]
// CoreIOT Configuration
#define CORE_IOT_SERVER "app.coreiot.io"
#define CORE_IOT_PORT 1883
#define DEVICE_TOKEN "YOUR_DEVICE_TOKEN_HERE" // Replaced with actual token

void publishTelemetry() {
    if (tb.connected()) {
        // Aggregate latest readings
        SensorData_t data;
        getSensorData(&data);
        
        // Publish to cloud
        tb.sendTelemetryData("temperature", data.temperature);
        tb.sendTelemetryData("humidity", data.humidity);
        tb.sendTelemetryData("light", data.light_level);
        
        Serial.println("[Cloud] Data published");
    }
}
\end{lstlisting}

\subsection{Integration Improvements}

Building upon the initial setup described in Chapter 1, we enhanced the cloud integration with several reliability features:

\begin{itemize}
    \item \textbf{Data Buffering:} Sensor readings are buffered in the FreeRTOS queue, ensuring that the latest valid reading is always available for publishing, even if the sensor task is temporarily busy.
    \item \textbf{Connection Resilience:} The system automatically handles MQTT broker disconnections without losing state. It attempts to reconnect to the cloud independently of the WiFi connection status.
    \item \textbf{Traffic Optimization:} Telemetry is published at a fixed interval (10 seconds) rather than on every sensor change, reducing bandwidth usage and complying with CoreIOT rate limits.
    \item \textbf{Bi-directional Control via RPC:} RPC callbacks (\texttt{setFanStatus}, \texttt{setLedEnabled}, \texttt{setWaterPumpStatus}) apply cloud commands directly to actuators and mirror state back with getter RPCs to keep dashboard toggles in sync.
    \item \textbf{Network Health Telemetry:} WiFi RSSI and derived quality percentage are published so the dashboard can surface link quality for troubleshooting.
\end{itemize}

This dual-interface approach—local Web Server for low-latency control and CoreIOT Cloud for historical analysis—provides a comprehensive user experience that combines the best aspects of edge and cloud computing.

\section{Data Logging and CSV Export}

To facilitate offline analysis and historical data tracking, we implemented a local data logging feature. The system periodically records sensor metrics and system states to a CSV (Comma-Separated Values) file stored in the ESP32's flash memory using the LittleFS file system.

\subsection{CSV File Structure}

The log file is named \texttt{data.csv} and contains the following columns:
\begin{itemize}
    \item \textbf{datetime:} Human-readable timestamp (YYYY-MM-DD HH:MM:SS)
    \item \textbf{timestamp:} Unix timestamp for machine processing
    \item \textbf{temperature, humidity:} Environmental readings
    \item \textbf{light\_level, moisture\_level:} Analog sensor values
    \item \textbf{flame\_detected, water\_pump:} Boolean status flags (0/1)
    \item \textbf{system\_state:} Textual representation of the system's operating mode
\end{itemize}

\subsection{System State Classification}

The \texttt{system\_state} column provides a high-level summary of the device's condition at the time of logging. This value is derived from the \texttt{SystemState\_t} enum and is exported as a string to ensure readability. Table~\ref{tab:csv_states} lists the possible values and their corresponding conditions.

\begin{table}[H]
\centering
\caption{System State values in the exported CSV file.}
\label{tab:csv_states}
\begin{tabular}{|l|l|l|}
\hline
\textbf{CSV Value} & \textbf{Enum Value} & \textbf{Condition} \\ \hline
Normal & STATE\_NORMAL (0) & Temperature 15--33\textdegree C, Humidity 50--70\% \\ \hline
Warning & STATE\_WARNING (1) & Temperature 33--40\textdegree C or Humidity out of ideal range \\ \hline
Critical & STATE\_CRITICAL (2) & Temperature $>$ 40\textdegree C or Humidity $<$ 30\% \\ \hline
Fire & STATE\_FIRE\_ALERT (3) & Flame sensor detected fire source \\ \hline
\end{tabular}
\end{table}

Users can download this file directly from the web interface for analysis in spreadsheet software like Microsoft Excel or Google Sheets.
